\section{Activity along long diagonal segments}\label{sec:segment-entropy}

The root probabilities constructed above detect variation within a plaque at
the intermediate resolutions.  We now use the arithmetic tube estimate to
show that this variation cannot disappear across a cut for a positive
proportion of starting points over a long diagonal segment.  A positive
average under $\mu_i$ alone could coexist with zero activity on a set of
starting points of positive measure.  If $\mu_i(E_i)\geq p>0$, the
normalized restrictions $\xi_i=\mu_i|_{E_i}/\mu_i(E_i)$ satisfy
$\xi_i\leq p^{-1}\mu_i$.  Proving the estimate for every such dominated
law excludes these positive-mass exceptional sets.  In
Section~\ref{sec:combination}, this stronger conclusion will give
almost-sure segment positivity for the limiting functions.

An \emph{oriented cut} is a nonempty proper subset
$J\subset\{1,\ldots,6\}$, with the roots directed from $J$ to its complement:
\[
 \mathcal E_J=\{\beta_{jk}:j\in J,\ k\notin J\}.
\]
The complementary subset gives the opposite orientation.  Fix a small
$\delta>0$.  For each such $J$, choose a regular $b_J\in L$ such that
\begin{equation}\label{eq:cut-direction}
 \beta(b_J)\geq1\quad(\beta\in\mathcal E_J),\qquad
 |(b_J)_k-(b_J)_j|\leq\delta
 \quad(j,k\text{ on the same side of }J).
\end{equation}
These choices can be made in a fixed bounded subset of $L$: put the two
sides at two distinct levels, perturb the entries within each side, and
subtract their common average.  There are only finitely many cuts.
Write $\eta=1/3$ for the exponential constant in
Proposition~\ref{prop:bowen}; it is independent of $\delta$.
Write $g_{\beta,i}=g_{jk,i}$ when $\beta=\beta_{jk}$, and put
\[
 A_{J,i}(x)=\sum_{\beta\in\mathcal E_J}g_{\beta,i}(x).
\]
Choose the common interior scale $R_i$ as in
\eqref{eq:wall-invariance}, leaving separated interior depths both below and
above it.  In particular, $T_{-,i}\prec R_i\prec T_{+,i}$.

\begin{proposition}[Activity on every long segment]\label{prop:segment-activity}
For $\delta$ sufficiently small, the preceding choices have the following
property.  For every oriented cut $J$, every $C<\infty$, every $\tau>0$,
and every sequence of probabilities $\xi_i$ on $X$ satisfying
$\xi_i\leq C\mu_i$,
\begin{equation}\label{eq:segment-activity}
 \liminf_{i\to\infty}\frac1{R_i}
 \int_0^{\tau R_i}\int_X
 A_{J,i}\bigl(xa(tb_J)\bigr)\,\dd\xi_i(x)\,\dd t>0.
\end{equation}
The same assertion holds after passage to any subsequence.
\end{proposition}

The proof compares two bounds for the entropy of orbit names.  The tube
estimate gives a positive entropy cost per unit time.  If the roots crossing
the cut are inactive, the plaque probabilities vary only in directions
whose expansion rates are at most $\delta$.  Their names then have a much
smaller coding cost.  This use of arithmetic separation to bound the entropy
of diagonal orbit names belongs to the method of
\cite{ELMV2009}; here the conditioning at intermediate resolutions and the
arbitrary dominated starting probabilities require the additional details
below.  We use Shannon entropy with natural logarithms, including its chain
rule and the elementary finite-alphabet error bound; see
\cite[Chapter~2]{CoverThomas2006}.

\subsection{Partitions and the cusp}

For a countable partition $\mathcal P$ and a probability $\nu$, write
\[
 H_\nu(\mathcal P)=-\sum_{P\in\mathcal P}\nu(P)\log\nu(P),
 \qquad 0\log0=0.
\]
Fix a regular $b\in L$ and put $T(x)=xa(b)$.  The notation
$\mathcal P_{[r,s]}=\bigvee_{t=r}^sT^{-t}\mathcal P$ records the symbols at
the integer times from $r$ through $s$.  Let $V$ be the pattern group of
the $b$-positive roots.

\begin{lemma}[A partition for consistent orbit names]\label{lem:entropy-partition}
Given a sufficiently small identity neighborhood $O\subset G$, there is a
countable Borel partition $\mathcal P$ of $X$ with the following properties.
\begin{enumerate}
 \item Every atom has $\mu$-null boundary and lies in an inner part of a
 relatively compact $VA$ plaque chart from the fixed stock.  The chart
 extends beyond the atom with a margin.
 \item If $x,y$ have the same symbols at times $0,\ldots,N$, there are
 relative lifts $h_t\in O$, $0\leq t\leq N$, such that
 \[
  ya(tb)=xa(tb)h_t,
  \qquad h_{t+1}=a(-b)h_ta(b).
 \]
 \item For every fixed $C$, the entropies $H_\nu(\mathcal P)$ are bounded
 uniformly for all sufficiently large $i$ and all probabilities
 $\nu\leq C\mu_i$.  Moreover, there are finite coarsenings
 $\mathcal P^{(q)}$, obtained by keeping finitely many atoms and combining
 the others into one atom, for which
 \begin{equation}\label{eq:partition-entropy-tails}
  \lim_{q\to\infty}\limsup_{i\to\infty}
  \sup_{\nu\leq C\mu_i}
  H_\nu(\mathcal P\mid\mathcal P^{(q)})=0.
 \end{equation}
\end{enumerate}
\end{lemma}

\begin{proof}
Let $m(x)$ denote the shortest-vector length.  Divide $X$ into layers on
which $m(x)$ is comparable to $e^{-k}$, for integers $k\geq0$, using a
first layer that also contains the thick part.  Layer boundaries may be
chosen $\mu$-null.  Successive minima and lattice basis reduction give a
representative matrix $g$ for each point of the $k$th layer such that
\[
 \|g\|+\|g^{-1}\|\leq e^{C_1(1+k)}.
\]
Consequently the injectivity radius for relative group coordinates on that
layer is at least $e^{-C_2(1+k)}$.  Indeed, if a nonidentity integral matrix
$\gamma$ gave a smaller return, conjugating the return by $g$ would give
$\|\gamma-1\|<1$, which is impossible.

We impose a quantitative margin condition on these charts.  Write $d_G$
for the fixed left-invariant Riemannian distance on $G$.  In a fixed
identity neighborhood of $G$, adapted product coordinates give a smooth
retraction $\pi$ onto $VA$, with $\pi(h)=h$ for $h\in VA$ and
\[
 d_G(h,\pi(h))\leq C_{\pi}\dist_G(h,VA),
 \qquad d_G(1,\pi(h))\leq C_{\pi}d_G(1,h).
\]
After shrinking the neighborhood, $1$ and $\pi(h)$ are joined inside
$VA$ by a path of length at most $C_{\pi}d_G(1,h)$; one may use the
local exponential coordinates of $VA$.

Choose outer $VA$ product charts in the $k$th layer of radius comparable
to $r_k=e^{-C_3(1+k)}$, with $C_3$ sufficiently large for injectivity,
and inner boxes whose relative-coordinate distance from the outer
boundary is at least a fixed multiple of $r_k$.  Use translates of a
fixed smooth group-coordinate chart, so that its local product and
metric comparison constants are uniform.  Cover the layer by still
smaller symbol boxes inside these inner boxes.  Require their
relative-group diameter to be at most $c_*r_k$, where $c_*>0$ is small
enough compared with both the inner-to-outer margin and the radius of
the domain of $\pi$.  In particular, whenever $x,y=xh$ are in one
symbol box, the whole path from $x$ to $x\pi(h)$ just described stays
inside its outer chart.  It follows that $x\pi(h)$ is on the same
chart plaque as $x$, rather than on another local sheet of that orbit.

The number of symbol boxes needed is at most $e^{C_4(1+k)}$.
Indeed, subdivide the bounded-matrix representatives at an exponentially
small mesh, and refine by the fixed factor $c_*$ when necessary.  These
operations preserve an exponential bound in $k$.  Passing from the cover
to a partition by successive differences does not increase the number
of pieces or their diameters.  The charts and their cuts can be chosen
from the stock with $\mu$-null boundaries: perturb centers within a small
fixed fraction of the covering radius and choose each coordinate width
in a fixed comparable interval, avoiding the at most countably many
widths that give a boundary of positive measure.  These choices preserve
the exponential bounds on sizes and on the number of boxes, as well as
the inner-chart margins.  Thus each partition atom remains
inside a chart with a margin, although it need not itself be a coordinate
rectangle.

The shortest-vector length changes by at most a fixed multiplicative factor
under $T$.  We can therefore choose the boxes so small that both a relative
lift in one symbol and its conjugate by $a(b)$ lie in an injective
neighborhood at the next symbol.  Uniqueness in that neighborhood identifies
the latter lift with the lift chosen there.  Induction gives the consistency
in assertion~2.  Decreasing all box sizes by a common fixed factor ensures
that every lift lies in the prescribed $O$.

By Proposition~\ref{prop:cusp}, a probability $\nu\leq C\mu_i$ assigns the
$k$th layer mass at most $C_5e^{-c k}$, uniformly for sufficiently large
$i$.  If that mass is $p_k$ and there are at most $M_k$ symbols in the
layer, their entropy is bounded by
\[
 p_k\log M_k-p_k\log p_k.
\]
Since $\log M_k=O(1+k)$ and $p_k\leq C_5e^{-ck}$, the sum of these bounds
converges uniformly, and its tail tends to zero uniformly.  Keeping all
symbols from the first finitely many layers proves both entropy assertions.
\end{proof}

\subsection{What an inactive cut implies inside a cell}

Fix $J$, put $b=b_J$, and let $V_s$ be the pattern of positive roots whose
endpoints lie on the same side of $J$.  Set $H_s=V_sA$.  The cross-root
group, denoted $Z_J$, is normal in $VA$, and each element of $VA$ has a
unique factorization
\[
 h=zd,\qquad z\in Z_J,\quad d\in H_s.
\]
In a product chart $x=\psi(q)zd$, the $H_s$ plaques are precisely the sets
with $(q,z)$ fixed.  Lemma~\ref{lem:va-support} says that at a limiting
state whose crossing-root indicators all vanish, the $VA$ probability
is supported on the $H_s$ plaque through the sampled point.

We need this support assertion for conditional probabilities of a dominated
measure, rather than only for $\mu_i$.  The following consequence records
the required change of measure.

\begin{lemma}\label{lem:inactive-cell-concentration}
Fix finitely many atoms $S$ of $\mathcal P$, with their $VA$ charts, and an
interior transverse depth $l_i\prec R_i$.  Suppose that probabilities
$\nu_i\leq C\mu_i$ satisfy $\int A_{J,i}\,\dd\nu_i\to0$.  Condition on
one of these symbols $S$ and on its transverse $q$ cell.  Draw two points
independently from the resulting conditional probability of $\nu_i$.
Then their $z$ coordinates differ by a quantity tending to zero in
probability, averaged with the probabilities of the conditioning events.

The conclusion also holds on average over any finite list of such
probabilities, with arbitrary probability weights, if their weighted mean
of $\int A_{J,i}\,\dd\nu_i$ tends to zero.
\end{lemma}

\begin{proof}
First sample $x$ under $\nu_i$, and draw the second plaque coordinate from
the $\mu_i$ probability conditioned only on its transverse cell in the
whole chart.  The depth-comparison rules allow this conditional probability
to be replaced, in the limit, by the chart probability in the state data.
Any joint limit of the state with $x$ sampled under $\nu_i$ is dominated by
$C$ times the usual state law.  Its crossing-root indicators vanish almost
surely, since their nonnegative sum has expectation tending to zero.
Lemma~\ref{lem:va-support} therefore puts the second point on the same
$H_s$ plaque as the sample.  Its $z$ coordinate equals the sampled $z$.
The transverse mesh tends to zero; continuous coordinate tests on the
fixed compact chart boxes now show that the prelimit $z$ difference tends
to zero in probability.  This argument uses the probability attached to
one chart and its statewise support.  It does not condition the sampled
point on the other plaque probabilities in the state.

To change the second draw to the stated $\nu_i$ conditional law, fix
$\varepsilon>0$ and discard events $E=S\cap\{q\text{ in a given cell}\}$
such that
\[
 \nu_i(E)<\varepsilon\mu_i(\{q\text{ in that cell}\}\cap\text{chart}).
\]
Their total $\nu_i$ mass is at most $\varepsilon$ times the number of
symbols under consideration.  On every remaining event the conditional
$\nu_i$ law has density at most $C/\varepsilon$ with respect to the
$\mu_i$ chart-cell law.  The preceding convergence thus remains valid for
the second draw.  Let $i\to\infty$ and then $\varepsilon\to0$.

For the averaged statement, first observe a uniform sequential consequence
of the first assertion: no sequence of dominated probabilities can have
indicator expectation tending to zero while the asserted pair error stays
bounded away from zero.  Let $a_i\to0$ be the weighted mean of the
indicator expectations in a list.  Discard members whose expectation is
greater than $\sqrt{a_i+i^{-1}}$.  Their total weight is at most
$a_i/\sqrt{a_i+i^{-1}}\to0$ by Markov's inequality.  On the remaining
members, the pair errors tend uniformly to zero by the sequential
consequence: otherwise a sequence of members violating it could be
selected.  This proves the averaged assertion.
\end{proof}

\subsection{The lower and upper entropy bounds}

\begin{proof}[Proof of Proposition~\ref{prop:segment-activity}]
Suppose, along a subsequence, that the expression in
\eqref{eq:segment-activity} tends to zero.  Choose even integers
$N_i\sim\tau R_i$ with $N_i+2<\tau R_i$.  Averaging a starting shift
$s\in[0,1]$ shows that we may replace $\xi_i$ by its translate through
some $s_i\in[0,1]$ so that
\begin{equation}\label{eq:vanishing-discrete-activity}
 \frac1{N_i+1}\sum_{t=0}^{N_i}
 \int A_{J,i}(T^t x)\,\dd\xi_i(x)\longrightarrow0.
\end{equation}
All these translated probabilities still satisfy $\xi_i\leq C\mu_i$.
More generally, every time-$t$ distribution
$\nu_{i,t}=(T^t)_*\xi_i$ has the same domination bound.

Choose a compact set $K$ with $\nu_{i,N_i/2}(K)\geq1-\varepsilon_0$ for
all sufficiently large $i$, where $\varepsilon_0>0$ will be small.  This
follows from the uniform cusp estimate and domination.  Take the
neighborhood $O$ in Lemma~\ref{lem:entropy-partition} small enough to use
Proposition~\ref{prop:bowen} uniformly on $K$.
If a name atom $B\in\mathcal P_{[0,N_i]}$ meets $T^{-N_i/2}K$, choose a
reference point in that intersection.  Consistency of the relative lifts
puts every midpoint of $B$ in the tube about its reference midpoint with
time $N_i/2$.  Since $T_{-,i}\prec R_i\prec T_{+,i}$, this time is in the
arithmetic range for all sufficiently large $i$.  Hence
\[
 \xi_i(B)\leq C e^{-\eta N_i/2}.
\]
Such atoms have total $\xi_i$ mass at least $1-\varepsilon_0$.
Integrating the negative logarithm of the atom mass gives, after decreasing
$\varepsilon_0$ and taking $i$ large,
\begin{equation}\label{eq:name-entropy-lower}
 \frac1{N_i}H_{\xi_i}(\mathcal P_{[0,N_i]})\geq\frac\eta3.
\end{equation}

We next derive the incompatible upper bound.  Fix a small target error
$\varepsilon>0$.  First choose a finite coarsening $\mathcal P^0$ so that
the conditional entropy in \eqref{eq:partition-entropy-tails} is as small
as required below.  Let $q_0$ be its alphabet size.  Choose also a finite
collection $\mathcal S$ of original symbols whose total complement has
arbitrarily small probability under every $\nu\leq C\mu_i$.  The symbols
in $\mathcal S$ need not be the symbols retained in $\mathcal P^0$.

The boundary of $\mathcal P^0$ is a finite union of compact $\mu$-null
sets.  We can choose a fixed distance accuracy $\rho>0$ so that its
$2\rho$ neighborhood has arbitrarily small probability under every
$\nu\leq C\mu_i$ for all sufficiently large $i$.  To justify this
uniformly, take decreasing closed neighborhoods of these boundaries;
their $\mu$ masses tend to zero, and the weak convergence upper bound
applies to each fixed neighborhood.  All subsequent orbit predictions
will be within distance $\rho$ of the actual point.  Incorrect symbols
can then occur only near these boundaries, or on the exceptions specified
below.

Choose transverse grids in the charts for $\mathcal S$ at an interior
depth $l_i\prec R_i$.  Let $M$ be a large integer, fixed while
$i\to\infty$.  Its size will be chosen after the charts and $\rho$;
the concentration tolerances used below may then depend on $M$.
Split the name into blocks of $M$ future symbols given the
past through the block's initial time.  We may omit an initial
$t_i=o(N_i)$ symbols and a final remainder of length less than $M$:
their entropy cost divided by $N_i$ tends to zero by the uniform
single-symbol entropy bound.

Here is why the extra transverse-cell label costs only a bounded amount
at every remaining block.  In a past-name atom at time $t$, the negative
root entries of its relative lift are $O(e^{-c_b t})$, where
$c_b=\min_{\beta(b)>0}\beta(b)>0$.  Triangular coordinates in a small
group neighborhood therefore put the lift at distance
$O(e^{-c_b t})$ from $VA$.  Let $x,y=xh$ be two current points in one
past-name atom, and use the retraction $\pi$ from the partition
construction.  The quantitative diameter condition ensures that
$x\pi(h)$ lies on the same plaque of the current outer chart as $x$.
Thus $q(x\pi(h))=q(x)$.  Lipschitz continuity of the transverse
coordinate on that chart gives
\[
 \begin{aligned}
 |q(y)-q(x)|
 &=|q(xh)-q(x\pi(h))|\\
 &\leq C\,d_G(h,\pi(h))
 \leq C'\dist_G(h,VA)
 \leq C''e^{-c_b t}.
 \end{aligned}
\]
The constants are uniform over the fixed finite collection of current
charts.  Choose
$t_i$ to be a sufficiently large fixed multiple of $l_i$; then
$e^{-c_b t}<2^{-l_i}$ for $t\geq t_i$, and $t_i=o(N_i)$.
Consequently the possible transverse coordinates meet only a bounded
number $B_0$ of neighboring cells.  Given the past, their label has
entropy at most $\log B_0$.  For current symbols outside $\mathcal S$
use a single dummy label.

Let $\mathcal Q_i$ record this cell label in the current symbol, with
the dummy label outside $\mathcal S$, and write
$\mathcal F_t=\mathcal P_{[0,t]}$.  After telling the cell label, we may
discard the rest of the past when upper-bounding the entropy of the next
coarsened block.  Precisely, the chain rule and monotonicity under
conditioning give
\begin{equation}\label{eq:coarse-block-entropy}
 \begin{aligned}
 &H_{\xi_i}\bigl((\mathcal P^0)_{[t+1,t+M]}\mid\mathcal F_t\bigr)\\
 &\quad\leq
 H_{\xi_i}\bigl(T^{-t}\mathcal Q_i\mid\mathcal F_t\bigr)
 +H_{\nu_{i,t}}\bigl((\mathcal P^0)_{[1,M]}
                     \mid\mathcal P\vee\mathcal Q_i\bigr)\\
 &\quad\leq\log B_0+
 H_{\nu_{i,t}}\bigl((\mathcal P^0)_{[1,M]}
                     \mid\mathcal P\vee\mathcal Q_i\bigr).
 \end{aligned}
\end{equation}
The conditioning in the last entropy is just the current symbol and
its cell.  By
\eqref{eq:vanishing-discrete-activity}, the mean crossing-root activity at
the block starts tends to zero.  Indeed their nonnegative sum is bounded
by the full sum in that equation, and normalizing by the number of blocks
loses at most the fixed factor $M$.  The averaged form of
Lemma~\ref{lem:inactive-cell-concentration} applies to their distributions.
For each symbol-cell event we may therefore choose a $z$ center so that,
apart from a probability tending to zero on average, $z$ lies within any
prescribed fixed tolerance of that center.  Such a center exists by
averaging the two-draw error and choosing one of the points as center.

The $d$ coordinates lie in a fixed compact part of $H_s$ for each of the
finitely many charts.  Every root occurring in $H_s$ has expansion rate
at most $\delta$, and $A$ commutes with the motion.  A mesh in these
coordinates, of accuracy a fixed multiple of $\rho e^{-\delta M}$,
therefore predicts the trajectory through the next $M$ steps to the
required distance accuracy, provided the $z$ and transverse coordinates
are sufficiently close to their centers.  The number of mesh points is
at most
\begin{equation}\label{eq:slow-direction-mesh}
 C_0\exp(C_6\delta M).
\end{equation}
Here $C_6$ depends only on the dimension, while $C_0$ depends on the fixed
charts and $\rho$, not on $M$.  This follows by using finitely many smooth
coordinate meshes on the compact $d$ ranges: conjugation expands every
entry of a group-relative error by at most $e^{\delta M}$.  Left
invariance of the metric transfers the resulting relative-group error
to an error on $X$.

The cross-root center tolerance may need to be exponentially small in
$M$, because cross-root directions expand faster.  This causes no
additional mesh labels: $M$ is fixed before applying the concentration
lemma, so any such positive tolerance is allowed.  The transverse width
also tends to zero for each fixed $M$.  Thus $C_0$ in
\eqref{eq:slow-direction-mesh} remains independent of $M$.

Encode the next $M$ symbols of $\mathcal P^0$ by a mesh label, with a
single additional label for an omitted chart or an outlier, and then by
corrections to the predicted symbols.  Let $e_{i,t}$ be the average
probability of an incorrect prediction over these $M$ positions, with
the starting point sampled under $\nu_{i,t}$.  Put
$h_2(e)=-e\log e-(1-e)\log(1-e)$, with the usual zero conventions.
The mesh and correction costs give
\begin{equation}\label{eq:mesh-correction-entropy}
 \begin{aligned}
 &H_{\nu_{i,t}}\bigl((\mathcal P^0)_{[1,M]}
                    \mid\mathcal P\vee\mathcal Q_i\bigr)\\
 &\quad\leq\log\bigl(C_0e^{C_6\delta M}+1\bigr)
       +M\bigl(h_2(e_{i,t})+e_{i,t}\log q_0\bigr).
 \end{aligned}
\end{equation}
For completeness, reveal the error indicator at each position and, on an
error, reveal the actual symbol; subadditivity and concavity of $h_2$
give this bound also when the error probabilities differ across positions
and blocks.  The chart omissions and boundary neighborhoods were chosen
with arbitrarily small probabilities.  The additional concentration
errors tend to zero on average.  Consequently the mean of $e_{i,t}$
over the block starts can be made as small as required; concavity of
$h_2$ gives the same small correction bound after averaging blocks.

Finally, the original block is recovered from its coarsened block at
the following cost:
\begin{equation}\label{eq:restore-block-entropy}
 \begin{aligned}
 &H_{\xi_i}\bigl(\mathcal P_{[t+1,t+M]}\mid\mathcal F_t\bigr)\\
 &\quad\leq
 H_{\xi_i}\bigl((\mathcal P^0)_{[t+1,t+M]}\mid\mathcal F_t\bigr)
 +\sum_{r=1}^M H_{\nu_{i,t+r}}(\mathcal P\mid\mathcal P^0).
 \end{aligned}
\end{equation}
Indeed, conditional on the entire coarsened block, revealing each
original symbol costs at most its one-step conditional entropy given
its own coarsened symbol.  Every summand on the right is controlled by
\eqref{eq:partition-entropy-tails}, because
$\nu_{i,t+r}\leq C\mu_i$.

Sum \eqref{eq:restore-block-entropy} over the blocks and use
\eqref{eq:coarse-block-entropy} and \eqref{eq:mesh-correction-entropy}.
The chain rule, together with the negligible initial and final pieces,
now gives
\begin{equation}\label{eq:name-entropy-upper}
 \limsup_{i\to\infty}\frac1{N_i}
 H_{\xi_i}(\mathcal P_{[0,N_i]})
 \leq C_6\delta+\varepsilon+
 \frac{\log B_0+\log(C_0+1)}{M}.
\end{equation}
All constants in the last numerator were fixed before $M$ was chosen.
Choose $\delta$ first so that $C_6\delta<\eta/12$, then choose the
truncation, charts and distance accuracy to make $\varepsilon<\eta/12$,
and finally choose $M$ so large that the last term is less than
$\eta/12$.  Letting $i\to\infty$ contradicts
\eqref{eq:name-entropy-lower}.  This proves the proposition, including its
subsequence formulation.
\end{proof}
